\documentclass{situationist-preprint}

\title{Charity and the Constraints on Translation}
\msid{LNG-06-b}
\msfamily{LNG-06}
\msdiscipline{Philosophy of Language}
\mstype{Consequence}
\mskeywords{radical translation; principle of charity; principle of humanity; Quine; Davidson; gavagai}

\begin{document}
\maketitle

\begin{abstract}
Discussions of radical translation usually treat the principle of charity as a basic constraint: a translation should maximize agreement between translator and native. LNG-06-a selects \enquote{rabbit} as the translation of \enquote{gavagai} by four constraints, none of which is charity in that form. We draw two consequences. First, charity is derivable from those constraints in the cases in which it is correct, and false in the cases in which it is not derivable: it holds where error would be unintelligible and fails where the native's circumstances make error intelligible. Second, charity is idle in the gavagai case, since the rival manuals attribute exactly the same agreement; the selection is made by uniformity and eligibility.
\end{abstract}

\section{Charity derived}

Quine held that assertions that are startlingly false on their face are likely to turn on hidden differences of language, and that the translator should prefer manuals that avoid them \citep{Quine1960}. Davidson made a stronger principle basic: an interpreter must so interpret as to make the speaker right, by the interpreter's lights, as far as possible \citep{Davidson1974}. Grandy proposed instead the principle of humanity, that the pattern of relations among beliefs, desires and the world attributed to the speaker be as similar to ours as possible \citep{Grandy1973}; LNG-06-a adopts a version of it as Principle~P4.

\begin{proposition}\label{p:charity}
Given Principles~P1 and~P4 of LNG-06-a, a manual that attributes to the native a false belief about her observable circumstances, which she has no intelligible reason to hold, is dispreferred to a manual that does not.
\end{proposition}

\begin{proof}[Argument]
A belief about observable circumstances that is false, and that the native's perceptual situation gives her no reason to hold, is unintelligible given her circumstances. Principle~P4 disprefers manuals that attribute unintelligible beliefs, other things equal. If both manuals satisfy Principle~P1, other things are equal with respect to stimulus fit, and the preference stands.
\end{proof}

Proposition~\ref{p:charity} is the content of charity in the cases where it is correct. It does not follow that agreement should be maximized. Where the native's circumstances make a false belief intelligible, because she lacks a piece of information we have, because she has been deceived, or because her perceptual position is poor, Principle~P4 does not disprefer the manual that attributes it; and if that manual does better on uniformity or eligibility (Principles~P2 and~P3), it is to be preferred, although it attributes less agreement. Maximizing charity would require the translator to reinterpret the native's words whenever she is intelligibly wrong, which is the error Grandy's principle was introduced to avoid. Charity, then, is not an independent constraint on translation; its correct applications are consequences of humanity.

\section{Charity in the gavagai case}

\begin{proposition}\label{p:idle}
The rabbit manual and its rivals, the undetached-rabbit-part manual and the rabbit-stage manual, attribute to the native the same distribution of true and false occasion sentences.
\end{proposition}

\begin{proof}[Argument]
By Principle~P1 of LNG-06-a, each manual maps the native's occasion sentences to sentences with the same stimulus meaning. \enquote{Lo, a rabbit}, \enquote{Lo, an undetached rabbit part} and \enquote{Lo, a rabbit stage} are true in the same observable circumstances. So each manual makes the native's assents true in the same circumstances.
\end{proof}

Charity, however formulated in terms of agreement, cannot choose between the manuals. LNG-06-a's selection of the rabbit manual does not depend on charity: Lemmas~3.1 and~3.2 there show that the rivals require non-uniform or gerrymandered translations of the native's predication and identity constructions, and Theorem~3.3 concludes that only the rabbit manual satisfies Principles~P2--P4. Humanity plays a confirming role, since a native interested in rabbit stages rather than rabbits would have interests unintelligible in her circumstances; but uniformity and eligibility settle the case before humanity is consulted.

\begin{objection}
Davidson argued that charity is not optional: without massive agreement, disagreement cannot be identified at all \citep{Davidson1974}.
\end{objection}

\begin{reply}
The present account grants that translation requires a large background of attributed true belief, but derives it. A native's beliefs about her observable surroundings are mostly intelligible only if mostly true, so humanity requires attributing mostly true observational beliefs. What it does not require is that agreement be maximized case by case, and Davidson's argument supports only the background, not the maximizing principle.
\end{reply}

\section{Conclusion}

Charity follows from the constraints that select \enquote{rabbit} where it is correct, fails where it diverges from humanity, and is idle in the gavagai case, which is settled by uniformity and eligibility.

\begin{thebibliography}{9}\small
\bibitem[Davidson(1974)]{Davidson1974} Davidson, D. (1974). On the very idea of a conceptual scheme. \emph{Proceedings and Addresses of the American Philosophical Association}, 47, 5--20.
\bibitem[Grandy(1973)]{Grandy1973} Grandy, R. (1973). Reference, meaning, and belief. \emph{Journal of Philosophy}, 70, 439--452.
\bibitem[Quine(1960)]{Quine1960} Quine, W.~V.~O. (1960). \emph{Word and Object}. MIT Press.
\end{thebibliography}

\end{document}
